\documentclass[11pt,a4paper]{article}

\usepackage[utf8]{inputenc}
\usepackage[T1]{fontenc}
\usepackage{amsmath,amssymb}
\usepackage{booktabs}
\usepackage{graphicx}
\usepackage{geometry}
\geometry{margin=2.2cm}
\usepackage{hyperref}
\hypersetup{colorlinks=true,linkcolor=blue,citecolor=blue,urlcolor=blue}

% ---------------------------------------------------------------------------
% IMPORTANT (author-internal note, remove before submission):
% This is an honest first draft built ONLY from frozen GPU results in
% reports/runs/. No statistic has been recomputed; every number below is
% transcribed from config.json / metrics.json / runs_index.json /
% engineer_status.md. R1 closed: main tables use *_repro_* runs at commit
% 458de18 (review round 2). kd_only has all 3 seeds; delta=0.5 locked and
% formally registered in configs/prereg/phase3_loss_config_v2.2.yaml.
% P2 fully closed: three seeds (seed0 delta_rand_audit / seed1 / seed2),
% commit 6784780, all mechanism_predictive, F1 pairwise-CI recomputation closed.
% P3/P4 not run. Per docs/protocol.md S10, contribution (1) plus the strong
% form of contribution (2) are claimed (scoped to the MVP equal-depth pair);
% contribution (3) is WITHHELD. Per decision tree (protocol S8) this manuscript
% remains a boundary / empirical study,
% NOT a high-level method paper.
% Figure placeholders are textual descriptions (no figure files exist yet).
% All citations marked [PLACEHOLDER] must be verified per protocol S9.
% ---------------------------------------------------------------------------

\title{Counterfactual Functional Fidelity in Cross-Model KV Handoff:\\
An Empirical Boundary Study}

\author{readkv project (author list to be finalized)}
\date{Draft --- generated from frozen results; not for distribution}

\begin{document}
\maketitle

% ===========================================================================
\begin{abstract}
Transferring a teacher model's key--value (KV) cache into a frozen student
promises generation with ``zero re-prefill'' of the source document. Standard
evaluation of such transfer asks whether the student still answers correctly.
We argue this is insufficient: a transferred cache can preserve ordinary
utility while distorting the model's \emph{counterfactual response}---the
change in its answer when a single fact in the document is edited. We define
an evaluation layer for this property, centered on the counterfactual error
$E_{\mathrm{CF}}$, and report a preregistered benchmark on Qwen3 model pairs
using real GPU runs (bf16, sdpa, strict zero re-prefill); the scoring code is
frozen at commit \texttt{458de18}, so this benchmark is reproducible from that
commit (Section~\ref{sec:models}). A measurement gate
(E0) shows implementation residuals at machine precision (logit cosine
$1.0$, KL $0$, measurement-error ceiling $2.2\times10^{-16}$), far below the
effect under study. On the equal-depth MVP pair (1.7B$\to$0.6B, three seeds),
the transferable Joint arm shows a counterfactual gap across the three seeds
reported here ($E_{\mathrm{CF}}\approx1.93$--$2.12$) whose bootstrap CI
excludes zero, with dual-side correctness near the native anchor and direction
consistency $0.90$--$0.96$. On the non-equal-depth main pair (4B$\to$1.7B,
three seeds), the gap is much larger ($E_{\mathrm{CF}}\approx5.37$--$6.08$)
but is accompanied by a severe collapse of transferred dual-side correctness
(from $0.65$--$0.78$ natively to $0.28$--$0.30$), and information controls are
mostly non-positive; we treat this pair as a broken-transfer failure boundary
and \emph{not} as positive evidence of counterfactual distortion. A matched
\texttt{kd\_only} baseline on the MVP pair (three seeds) narrows but does not
eliminate the affine gap: $E_{\mathrm{CF}}$ $0.735\;[0.552,0.934]$,
$1.034\;[0.766,1.328]$, and $1.539\;[1.148,1.959]$ versus affine
$1.968$/$1.934$/$2.123$, with all CI lower bounds above zero. Matched-budget
confirmation remains a caveat. On the MVP equal-depth pair (1.7B$\to$0.6B),
three seeds of held-out mechanism analysis (runs
\texttt{p2\_gpu\_mvp\_seed0\_delta\_rand\_audit}, \texttt{p2\_gpu\_mvp\_seed1},
\texttt{p2\_gpu\_mvp\_seed2}; commit \texttt{6784780}) yield
\texttt{mechanism\_predictive} in every seed. K/V four-arm errors are Self 0,
K-only 0.582, V-only 1.626, Joint 1.777 (only Joint deployable); a local
routing/consumption decomposition over 17,920 units has residual
$2.96\times10^{-13}$; and units chosen on 8 training pairs repair held-out
(40 pairs) cases more often than random units, with paired differences
$0.275\;[0.150,0.425]$, $0.275\;[0.125,0.425]$, and $0.350\;[0.200,0.500]$
(all lower bounds $>0$) and direction consistency $0.725$--$0.750$. We
\emph{localize the gap to receiver-side Value-consumption pathways under
controlled cache interventions}, with Key routing as an accompanying
attention-redistribution effect; this is a local, fixed-query decomposition,
not a global causal account, and the claim is confined to the MVP equal-depth
pair. Core units reproduce across seeds with one marginal unit drifting. With
evaluation-layer and mechanism evidence, but no method evidence, we present a
\emph{boundary/empirical study}: we identify a gap for the affine/ridge mapper
tested here, document the conditions under which the current transfer fails,
and report the held-out mechanism evidence. Whether other strong behavioral
baselines (\texttt{consumer}, \texttt{score}, ridge) retain the gap is left to
future work. We make \emph{no} claim of a causal or exclusive mechanism, of
behavior across model pairs, or of having developed a method that closes the
gap.
\end{abstract}

% ===========================================================================
\section{Introduction}
\label{sec:intro}

Cross-model KV handoff maps a teacher's document cache into a receiver
student's cache format so that the student can answer queries about a
document without ever re-processing that document (\emph{strict zero
re-prefill}). The natural question in prior work is whether the student still
produces the correct answer. We ask a sharper question: does the student
preserve the \emph{functional dependence of its answer on the document}? If a
target fact in the document is counterfactually edited, does the transferred
student's response change in the same direction and by a comparable magnitude
as the native student's response? We call this property \emph{counterfactual
functional fidelity} and measure it with the counterfactual error
$E_{\mathrm{CF}}$ (Section~\ref{sec:metrics}).

This question sits at the intersection of two lines of work (Figure~\ref{fig:f1}):
(i)~cross-model KV transfer, which optimizes standard utility and distribution
matching but does not score counterfactual response change; and (ii)~within-model
content/counterfactual auditing, which studies whether a model's output is
causally driven by document content but does not address transfer into a
frozen receiver. Existing transfer work therefore leaves open exactly the
property we evaluate, while existing auditing work does not cover the
zero-re-prefill setting. Our evaluation layer connects the two by measuring
\emph{counterfactual response fidelity after transfer} under a preregistered
protocol.

\paragraph{Contributions (strictly scoped).}
Following the contribution gate in the project protocol (protocol~\S10), we
claim exactly \emph{one} evidence-backed contribution:

\begin{enumerate}
\item \textbf{Evaluation layer (claimed, scope-limited).} For the
      affine/ridge mapper tested here, we identify a gap between standard
      transfer utility and counterfactual response fidelity in cross-model KV
      handoff, together with a preregistered \emph{reproducible} paired
      counterfactual benchmark (scoring code frozen at commit \texttt{458de18}),
      a machine-precision measurement gate, and a joint reporting discipline
      ($E_{\mathrm{CF}}$ reported with dual-side correctness, direction
      consistency, coverage, and information-control gaps). A matched
      \texttt{kd\_only} baseline (three seeds) narrows this gap without
      eliminating it ($E_{\mathrm{CF}}=0.735\;[0.552,0.934]$,
      $1.034\;[0.766,1.328]$, and $1.539\;[1.148,1.959]$ versus affine
      $1.968$/$1.934$/$2.123$; all CI lower bounds above zero). This claim is
      explicitly limited to
      the mapper studied: the other strong behavioral baselines
      (\texttt{consumer}, \texttt{score}, ridge) are not evaluated here and are
      left to future work. The three-seed \texttt{kd\_only} evidence shows a
      persistent gap, but its unequal training budget
      (Section~\ref{sec:pending}) and the absence of a trained CF-RCT method
      preclude any claim of a method-level gain at this stage.
\item \textbf{Mechanism layer.} On the MVP equal-depth pair (1.7B$\to$0.6B),
      three seeds (seed 0/1/2), 40 held-out pairs each, we
      \emph{localize the gap to receiver-side Value-consumption pathways under
      controlled cache interventions}, with Key routing as an accompanying
      attention-redistribution effect. Units selected on 8 training pairs
      repair held-out cases more often than random units, with paired
      differences $0.275\;[0.150,0.425]$, $0.275\;[0.125,0.425]$, and
      $0.350\;[0.200,0.500]$ (all lower bounds $>0$) and direction consistency
      $0.725$--$0.750$. Four core units reproduce across all three seeds and
      the fifth drifts marginally. This is a local, fixed-query decomposition,
      \emph{not} a global causal account; the three E2 criteria are two
      statistical signals plus a majority-threshold restatement, not three
      independent lines of evidence. The claim is confined to the MVP
      equal-depth pair; it makes no causal, exclusive, or cross-model-pair
      assertion.
\end{enumerate}

We do \emph{not} claim the following contribution at this stage, because the
required evidence has not been produced:
\begin{itemize}
\item \emph{(withheld) method layer:} the mechanism-constrained-mapper claim
      (contribution template~(3) in protocol~\S10), which would require a
      trained method evaluated against matched behavioral baselines on the
      sealed test.
\end{itemize}
The P1 directional clue (V-only $\approx$ Joint) was subsequently confirmed
and extended by the P2 intervention and held-out analysis reported in
Section~\ref{sec:results}. The proposed CF-RCT method
(Section~\ref{sec:method}) is described as a \emph{design pending validation};
it has not been trained or evaluated. Accordingly, this manuscript follows the
decision tree (protocol~\S8) and is positioned as a boundary/empirical study,
not a high-level method paper. All data used here are synthetic templates; we
therefore make no claim about cross-family or real-world transfer, and mark
this as an applicability boundary.

\begin{figure}[t]
\centering
\fbox{\parbox{0.88\linewidth}{\centering
\textbf{Figure~1 (placeholder; to be generated from frozen results).}\\
Two prior-work lines $\to$ the third scientific question.\\
Left: \emph{cross-model KV transfer} (standard utility / distribution matching).\\
Right: \emph{within-model content \& counterfactual auditing}.\\
Center/foreground: \emph{this work}---counterfactual response fidelity of a
frozen receiver under strict zero re-prefill, evaluated with paired
counterfactual documents and a preregistered control suite.\\
Annotate that prior transfer work does not score counterfactual response change,
and prior auditing does not cover zero-re-prefill transfer.}}
\caption{Relation of the two established research lines to the scientific
question of this work. \emph{Textual placeholder; figure file not yet generated.}}
\label{fig:f1}
\end{figure}

% ===========================================================================
\section{Related Work}
\label{sec:related}

\paragraph{Cross-model KV transfer.} Methods that distill or map a teacher's
caches into a receiver's cache format (e.g., KV-Lingo-style behavioral
distillation and CacheLink-style cache translation)~\cite{kvlingo,cachelink}
evaluate whether the receiver retains ordinary task accuracy or matches the
teacher's token distributions. They do not score whether the receiver's
\emph{response difference} under a counterfactual document edit is preserved.
We treat these as strong baselines to be matched on data, forward budget,
parameters, and tuning budget (Section~\ref{sec:setup}); we do not assume their
reported numbers transfer to our task, and where no runnable implementation is
available we report the deviation explicitly, and we do not assume the
reported numbers carry over to our task.

\paragraph{Within-model auditing and counterfactual evaluation.} Content-
and counterfactual-dependence audits test whether a model's output is causally
attributable to document content. We reuse the principle of comparing against
no-information and mismatched-content controls, but we do not claim a novel
causal-audit metric; the counterfactual error and its paired construction are
an evaluation design for the transfer setting.

\paragraph{Scope boundaries.} Four categories are kept distinct throughout:
(1)~cross-model transfer (our main line); (2)~within-model incremental cache
editing/repair (not studied); (3)~no-information message adaptation (used only
as a control); and (4)~content-dependence auditing (adopted principle, not a
claimed contribution). We further do not claim that representation error or
$R^2$ equals task quality, that single-document attention error explains a
counterfactual effect, that offline Oracle diagnostics are deployable, or that
identity/cache round-trip auditing is itself a contribution.
 All citations in
this draft are placeholders pending version/date/protocol verification under
protocol~\S9.

% ===========================================================================
\section{Method}
\label{sec:method}

This section is organized in two parts. Sections~\ref{sec:proposed}--%
\ref{sec:huber} describe the CF-RCT design as a \emph{proposed method not yet
validated}. Section~\ref{sec:zerore} defines the evaluation pathway that
\emph{was} realized in the reported runs.

\subsection{Setup and notation}
\label{sec:proposed}

A teacher $T$ prefill-computes a cache $C_T(d)$ of document $d$; a trainable
mapper $g_\theta$ maps it into the receiver $S$'s cache format
$\hat{C}_S(d)=g_\theta(C_T(d))$. The receiver decodes query $q$ starting from
$\hat{C}_S(d)$ and never re-processes $d$. The receiver's own cache
$C_S(d)$ is an \emph{offline Oracle}: it may be used for training targets,
evaluation anchors, and mechanism intervention, but never inside the
deployable inference path. $S$'s main weights are frozen; $g_\theta$ is the
only trainable object. The two documents of a counterfactual pair, $d_A$ and
$d_B$, are prefilled \emph{independently}; cache bytes must never be tampered
with to synthesize the other document.

\subsection{Strict zero re-prefill and the arm boundary}
\label{sec:zerore}

We distinguish four arms by the provenance of Key and Value (Figure~\ref{fig:f2},
Table~\ref{tab:arms}). Only \emph{Joint} is deployable under strict zero
re-prefill: the student processes no document token and does not read
$C_S(d)$ at inference. A forward hook asserts that the number of student
forward passes equals one plus the decode steps and that no document token
enters the student; a per-run \texttt{zero\_re\_prefill} flag records this.

\begin{figure}[t]
\centering
\fbox{\parbox{0.88\linewidth}{\centering
\textbf{Figure~2 (placeholder).}\\
Strict zero re-prefill flow and diagnostic boundary.\\
Teacher independently prefills $d_A$, $d_B$ $\to$ $C_T$; mapper
$\hat{C}=g_\theta(C_T)$; student decodes $q$ from $\hat{C}$ only
(\emph{Joint}, deployable).\\
Dashed/offline box (NOT in the deployable graph): student native $C_S(d)$
(\emph{Self}), identity injection (\emph{Identity}), and layer/head/token
\emph{Oracle} state grafting.\\
K-only / V-only mix one mapped and one native half and are marked
\texttt{deployable=false}.\\
Label the only honest comparison: Joint vs.\ matched baselines; all others are
offline diagnostics or anchors.}}
\caption{Strict zero re-prefill inference and the Self / Identity / Joint /
Oracle boundary. \emph{Textual placeholder; figure file not yet generated.}}
\label{fig:f2}
\end{figure}

\begin{table}[t]
\centering
\caption{Arm definitions and deployability. Only Joint satisfies strict zero
re-prefill; other arms are offline diagnostics/anchors.}
\label{tab:arms}
\begin{tabular}{llll}
\toprule
Arm & Key & Value & Deployable \\
\midrule
Self    & student $K_S$   & student $V_S$    & no (anchor) \\
Identity& student $K_S$   & student $V_S$    & no (anchor) \\
K-only  & mapped $\hat{K}_S$ & student $V_S$  & no \\
V-only  & student $K_S$   & mapped $\hat{V}_S$ & no \\
Joint   & mapped $\hat{K}_S$ & mapped $\hat{V}_S$ & \textbf{yes} \\
Oracle  & student-native state grafted into mapped cache & --- & no \\
\bottomrule
\end{tabular}
\end{table}

\subsection{Proposed CF-RCT algorithm (design, not validated)}
\label{sec:cf-rct}

We outline a family \texttt{CF-RCT} (`counterfactual response-constrained
translation') with two instances, \texttt{CF-RCT-lite} and
\texttt{CF-RCT-mech}, sharing a training skeleton. The mapper is a
per-(layer, channel$\in$\{K,V\}, head) affine map with an optional low-rank
residual $W=W_0+AB^\top$ (rank $0$ is the pure-affine special case). For non-equal-depth
pairs (4B$\to$1.7B, $36\to28$ layers) a proportional layer map is used rather
than forcing same-layer alignment. Initialization is a closed-form ridge
regression from teacher to student caches on a calibration set; computation is
done in de-RoPE'd content space and re-RoPE'd before injection, with bf16 +
sdpa in every arm.

For each calibration pair, two student forward passes under the mapped caches
produce $D_{\mathrm{trans}}$, while a \texttt{no\_grad} pass under the native
caches produces the offline anchor $D_{\mathrm{self}}$. The training objective is
\begin{equation}
L \;=\; \lambda_{\mathrm{KD}} L_{\mathrm{KD}}
      + \lambda_{\mathrm{CF}} L_{\mathrm{CF}}
      + \lambda_{\mathrm{local}} L_{\mathrm{local}}
      + \lambda_{\mathrm{reg}} L_{\mathrm{reg}},
\end{equation}
where $L_{\mathrm{KD}}$ is the per-document forward KL between mapped and
native candidate distributions (the strong behavioral baseline
\texttt{kd\_only} corresponds to $\lambda_{\mathrm{CF}}=0$);
$L_{\mathrm{CF}}=\mathrm{mean}_{i}[m_i\,\mathrm{Huber}_{\delta}
(D_{\mathrm{trans}}^{i}-D_{\mathrm{self}}^{i})]$ is a paired counterfactual
residual constrained only on the ``Self-valid'' subset
$m_i=\mathbf{1}[\,|D_{\mathrm{self}}|\ge 0.1\,]$; $L_{\mathrm{local}}$ is a
routing/consumption term used only by \texttt{mech} and only in training; and
$L_{\mathrm{reg}}$ penalizes the low-rank residual. `lite' sets
$\lambda_{\mathrm{local}}=0,\lambda_{\mathrm{CF}}>0$; `mech' additionally sets
$\lambda_{\mathrm{local}}>0$. A key design guard is that $L_{\mathrm{CF}}$
is computed only where the native signal is identifiable: forcing sensitivity
where $D_{\mathrm{self}}\approx 0$ can induce a degenerate ``constant output''
solution (Section~\ref{sec:metrics}).

\paragraph{Mechanism diagnostics (executed at P2).} A routing/consumption
decomposition at fixed receiver query $Q_S$ is the algebraic identity
$\hat O-O_S=(\hat A-A_S)V_S W_O+\hat A(\hat V-V_S)W_O$, whose two terms are
\emph{routing} and \emph{consumption}. This is a \emph{local, fixed-query
identity}, not a global causal decomposition of task loss; $d_A$ and $d_B$ are
diagnosed separately and their differences compared. Layer/head/token Oracle
grafting ranked units by repair gain on 8 training pairs and tested them on 40
held-out validation pairs with no overlap and no use of the sealed test;
Oracle states never enter the Joint inference graph. At P2 (three runs
\texttt{p2\_gpu\_mvp\_seed0\_delta\_rand\_audit}, \texttt{p2\_gpu\_mvp\_seed1},
\texttt{p2\_gpu\_mvp\_seed2}; commit \texttt{6784780}) the decomposition
residual over 17,920 units was $2.96\times10^{-13}$, and the held-out repair
hit rate exceeded the random baseline in every seed (paired differences
$0.275\;[0.150,0.425]$, $0.275\;[0.125,0.425]$, $0.350\;[0.200,0.500]$;
direction consistency $0.725$--$0.750$). Details are in
Section~\ref{sec:results}.

\subsection{Hyperparameter freezing status}
\label{sec:huber}

The Huber threshold $\delta$ (nominal $1.0$, logit units), label balance
(A/B $1{:}1$, candidate order 50/50, F5 negative control fixed, $L_{\mathrm{CF}}$
mask on Self-valid subset) and loss-weight selection rules are frozen by a
dedicated preregistration file (protocol-v2.1). The validation grid
$\delta\in\{0.5,1.0,2.0\}$ has now been completed, and the PI locked the final
value at $\delta=0.5$ under the hard-guardrail reading of ``not worse than
\texttt{kd\_only}'' (Section~\ref{sec:adverse}); $\delta=1.0$ and $\delta=2.0$
were rejected for violating the dual/flip guardrail. The lock is on validation
data only, formally registered in \texttt{configs/prereg/phase3\_loss\_config\_v2.2.yaml}
(protocol-v2.2), and the value must not be adjusted during the sealed test.
Still, no
\texttt{CF-RCT} model was trained and the sealed test was never run, so we
report no method result; the grid evidence in fact warns that no $\delta$
beat \texttt{kd\_only} on the validation $E_{\mathrm{CF}}$.

% ===========================================================================
\section{Experimental Setup}
\label{sec:setup}

\subsection{Counterfactual data construction: slices F1--F6}
\label{sec:data}

Each evaluation unit is $u=(q,d_A,d_B,a,b,m)$: a query and a pair of documents
differing in one targeted fact, with candidates A/B held fixed within the pair
and their order randomly swapped. Six slices are designed to exclude distinct
confounds (Figure~\ref{fig:f3}; protocol~\S3.2):
\begin{itemize}
\item \textbf{F1} direct fact substitution (who/when/how-many/where): holds
      position and overall context nearly fixed;
\item \textbf{F2} reasoning-chain change (alias/membership/comparison) with
      the answer in a different span: excludes answer-string shortcuts;
\item \textbf{F3} relation direction / negation reversal: excludes lexical
      overlap with unchanged logical relation;
\item \textbf{F4} distractors and coreference: excludes selection by entity
      frequency or position;
\item \textbf{F5} surface equivalence (paraphrase/word order/format): a
      \emph{negative control} that should \emph{not} flip the answer; excluded
      from the ``must-flip'' denominator;
\item \textbf{F6} length/position control (matched token lengths; separate
      mismatched and position-shifted subsets): excludes attributing RoPE /
      position effects to fact infidelity.
\end{itemize}
Splits are isolated by entity, source, and template cluster; train/val/test
counts are 120/60/300 with 40 probe pairs, and the train--eval template overlap
is empty in the runs (data-quality gate `pass'). The reported P1 runs evaluate
on the validation split (60 pairs); the sealed test is reserved for the
(unrun) method stage.

\begin{figure}[t]
\centering
\fbox{\parbox{0.88\linewidth}{\centering
\textbf{Figure~3 (placeholder).}\\
Paired counterfactual document construction and positive/negative controls.\\
Show $d_A$ vs.\ $d_B$ differing in one edited span; candidates A/B order
randomly swapped; F1--F6 examples with F5 flagged as non-flip negative control
and F6 as length/position control.}}
\caption{Construction of counterfactual document pairs and controls.
\emph{Textual placeholder; figure file not yet generated.}}
\label{fig:f3}
\end{figure}

\subsection{Model pairs}
\label{sec:models}

Two real-GPU pairs were run: (a) \emph{MVP}, equal-depth Qwen3-1.7B (teacher)
$\to$ Qwen3-0.6B (student), $28\to28$ identity layer map; and (b) \emph{main
validation}, non-equal-depth Qwen3-4B $\to$ Qwen3-1.7B, $36\to28$ proportional
layer map. All runs use device CUDA, dtype \texttt{bfloat16}, attention kernel
\texttt{sdpa}, protocol version 2.0, and \texttt{zero\_re\_prefill=true}.
Three seeds ($0,1,2$) were run for each P1 pair. The 8B offline stress pair was
not run.%
\footnote{Data-integrity note (model pair): the \texttt{model\_pair} string in
the main runs' \texttt{config.json}/\texttt{runs\_index.json} reads
\texttt{qwen3-1.7b-to-0.6b}, but the actual \texttt{teacher}/\texttt{student}
paths (\texttt{Qwen3-4B}/\texttt{Qwen3-1.7B}) and the $36\to28$ proportional
layer map confirm these are Qwen3-4B$\to$1.7B runs. We therefore report the
corrected label \texttt{qwen3-4b-to-1.7b}. The metadata erratum is recorded in
\texttt{reports/runs\_index\_errata.json} (\texttt{model\_pair\_corrected});
the frozen \texttt{scores.jsonl} were not modified, and no computed metric is
affected.}}
\paragraph{Provenance note.} Every run used for the main numeric tables was
produced from commit \texttt{458de18}, which contains all scoring and
implementation code; the benchmark is therefore reproducible from that commit.
The authoritative runs are \texttt{e0\_gpu\_mvp\_repro\_seed0} and
\texttt{p1\_gpu\_mvp\_repro\_seed\{0,1,2\}} (plus the \texttt{kd\_only} run
described in Section~\ref{sec:pending}); each records
\texttt{git\_hash=458de18}, \texttt{bootstrap\_seed=seed},
\texttt{zero\_re\_prefill=true}, \texttt{dtype=bfloat16}, \texttt{kernel=sdpa},
and \texttt{protocol\_version=2.0}. The earlier eight runs carry
\texttt{git\_hash=53e9bc5} (a tree without the scoring code) and are marked
\texttt{superseded / provenance-unverified} in
\texttt{reports/runs\_index\_errata.json}; their frozen artifacts were not
modified, and they are not used for any main-table number. The
non-equal-depth pair has no \texttt{458de18} rerun, so its results
(Section~\ref{sec:main}) are retained only as a provenance-limited
broken-transfer record, not as authoritative numeric evidence.

\subsection{Metrics: $E_{\mathrm{CF}}$ and mandatory joint reporting}
\label{sec:metrics}

For a candidate pair, define $g(C,q)=\log P(A\mid C,q)-\log P(B\mid C,q)$ at a
single shared generation position. Then
\begin{align}
D_{\mathrm{self}} &= g(C_S(d_A),q)-g(C_S(d_B),q),\\
D_{\mathrm{trans}} &= g(\hat C_S(d_A),q)-g(\hat C_S(d_B),q),\\
E_{\mathrm{CF}} &= |D_{\mathrm{trans}}-D_{\mathrm{self}}|,\qquad
\mathrm{normalized}\ E_{\mathrm{CF}}=\frac{|D_{\mathrm{trans}}-D_{\mathrm{self}}|}
{|D_{\mathrm{self}}|+\tau},\quad \tau=1.0 .
\end{align}
The key interpretation guard: $E_{\mathrm{CF}}$ alone is not a quality measure.
When $D_{\mathrm{self}}\approx0$, a constant output can achieve small
$E_{\mathrm{CF}}$. Therefore $E_{\mathrm{CF}}$ is \emph{always} jointly reported
with (i)~dual-side per-document success \texttt{dual\_correct}, (ii)~direction
consistency \texttt{flip\_consistency} on the Self-valid subset
$|D_{\mathrm{self}}|\ge0.1$, (iii)~coverage \texttt{coverage} of the native
evaluable subset, and (iv)~the gap to no-information controls
\texttt{info\_gap}/\texttt{info\_attribution}. A decrease in $E_{\mathrm{CF}}$
without these joint signals is not counted as an improvement.

\paragraph{Information-attribution reporting note.} In the current pipeline the
control groups are emitted as separate arms, so the per-pair
\texttt{info\_gap} field is unset; the aggregate equivalent
\texttt{info\_attribution} (control-group $E_{\mathrm{CF}}$ minus Joint
$E_{\mathrm{CF}}$) is reported instead. This deviation is recorded, not used
to alter any frozen criterion. Positive values mean the correct cache carries
identifiable information beyond the controls.

\subsection{Primary endpoint and statistical discipline}
\label{sec:stats}

The primary endpoint is defined for the \emph{method} stage: the paired
$E_{\mathrm{CF}}$ difference on the sealed test between a new method and the
strongest behavioral baseline (\texttt{primary\_baseline=kd\_only}), with a
template/source-cluster paired bootstrap (10{,}000 resamples, 95\% CI). An
improvement is declared only if the CI \emph{upper bound} is $<0$; $\ge3$
training seeds are required and multiple comparisons use Holm. A result with
$P>0.05$ is not treated as evidence of equivalence. The P1 runs reported here
are the \emph{benchmark} stage: they use no new method, so the paired endpoint
is degenerate (point $0$) and \emph{no improvement claim is made}; their CIs
quantify the gap of Joint above the Self anchor (which has $E_{\mathrm{CF}}=0$).

\subsection{Baselines and fair alignment}
\label{sec:baselines}

Design requirement: \texttt{affine}, \texttt{consumer}, \texttt{kd\_only},
\texttt{score}, plus ridge, CE-only, contrastive, mismatched-document and
receiver-only-adapted controls, all on the same calibration pairs, dtype/kernel,
forward budget, trainable parameter count, training steps, and tuning budget,
with ordinary QA and counterfactual QA reported together. In the P1 runs the
mapped arm is an affine mapper with \emph{zero} trainable parameters (closed-form
ridge initialization), so the headline comparison here is Joint vs.\ Self /
Identity and against mismatched / random / zero information controls; the
full matched-baseline comparison belongs to the unrun method stage.

% ===========================================================================
\section{Results}
\label{sec:results}

We report only real GPU results from \texttt{reports/runs/}. No mock number is
used. All values are transcribed from frozen \texttt{config.json}/\texttt{metrics.json}
and the engineering status report; no statistic was recomputed for this draft.
The main numeric tables use the \texttt{458de18} authoritative runs listed in
Section~\ref{sec:models}; the superseded runs appear only in the provenance
footnotes.%
\footnote{Bootstrap CI endpoints depend on the bootstrap seed; point estimates
are exact. The scoring pipeline and the audit script now both read
\texttt{bootstrap\_seed} from each run config, and the endpoint differences
between the two were resolved by reviewer item R5 (closed in review round 2).}

\subsection{E0: implementation gate (P0)}
\label{sec:e0}

E0 isolates pure implementation residual by using one model as both capture and
inject (identity mapper) on 40 probe pairs. Both model pairs pass:
\begin{table}[t]
\centering
\caption{E0 implementation gate, authoritative run \texttt{e0\_gpu\_mvp\_repro\_seed0}
(commit \texttt{458de18}, real GPU). All thresholds met; the measurement-error
ceiling is orders of magnitude below the effect under study. ``cosine'' refers
to first-token logit cosine on the candidate pair.}
\label{tab:e0}
\begin{tabular}{lcc}
\toprule
Check (threshold) & MVP 1.7B$\to$0.6B (repro) & Pass \\
\midrule
Self/Identity cosine ($\ge0.9998$) & $1.000000$ & yes \\
First-token KL ($\le0.001$) & $0.000000$ & yes \\
Argmax/EM agreement ($\ge0.929$) & $1.000$ (80/80) & yes \\
Round-trip bit-identical & True (80/80) & yes \\
zero\_re\_prefill assert & 80/80 True & yes \\
Measurement-error ceiling & $2.2\times10^{-16}$ & yes \\
Effect estimate / tolerance (ceiling $<$ effect/3) & $20.86$ / $6.95$ & yes \\
\bottomrule
\end{tabular}
\end{table}
A parallel E0 check on the 1.7B student for the non-equal-depth pair was
performed in an earlier run (\texttt{e0\_gpu\_main\_seed0}) that is now
superseded and has no \texttt{458de18} equivalent; its values are not used as
evidence here.
A supplementary diagnostic measures the full-vocabulary residual of a one-shot
whole-document prefill against the cached-plus-query two-step path
(min cosine $0.9993$, 47/80 below $0.9998$); this comes from the sdpa kernel
taking different paths for whole-prefill vs.\ cached+query, not an
implementation error, and is excluded from the E0 anchors (the injection path
is bit-identical to the native two-step path, cosine $1.0$).

\subsection{P1 benchmark, MVP pair (equal depth, 1.7B$\to$0.6B)}
\label{sec:mvp}

\begin{table}[t]
\centering
\caption{P1 MVP pair, authoritative runs \texttt{p1\_gpu\_mvp\_repro\_seed\{0,1,2\}}
(commit \texttt{458de18}), three seeds, validation split (60 pairs). Joint
$E_{\mathrm{CF}}$ with 95\% cluster-paired bootstrap CI; joint dual correctness,
flip consistency, coverage, and normalized $E_{\mathrm{CF}}$ reported jointly.
Self/Identity have $E_{\mathrm{CF}}=0$ by construction. The superseded
non-repro runs give identical values but are not cited here.}
\label{tab:mvp}
\begin{tabular}{cccccc}
\toprule
seed & Joint $E_{\mathrm{CF}}$ [95\% CI] & dual\_correct & flip\_consistency & coverage & norm.\ $E_{\mathrm{CF}}$ \\
\midrule
0 & $1.968\;[1.614,2.324]$ & $0.667$ & $0.938$ & $0.800$ & $0.334$ \\
1 & $1.934\;[1.569,2.311]$ & $0.650$ & $0.898$ & $0.817$ & $0.290$ \\
2 & $2.123\;[1.695,2.571]$ & $0.617$ & $0.958$ & $0.800$ & $0.322$ \\
\bottomrule
\end{tabular}
\end{table}

Across the three seeds reported here, the Joint counterfactual gap has a CI
whose lower bound lies above zero:
transfer is \emph{not} counterfactually faithful. The Self anchor's own dual
correctness is $0.683/0.600/0.650$ (the synthetic task is not trivially solved
by the native model), and Joint stays close to it. Direction consistency is
high ($0.90$--$0.96$), and coverage $\approx0.80$. Information controls on seed 0
(control $E_{\mathrm{CF}}-$ Joint $E_{\mathrm{CF}}$: mismatch $+4.134$,
random $+2.249$, zero $+4.191$) are positive: the correct cache carries
identifiable information, so the observed gap is not a ``no-information''
artifact. The K/V four arms (seed 0: K-only $0.630$, V-only $1.851$,
Joint $1.968$; similar in seeds 1 and 2) show V-only close to Joint---the
\emph{directional clue} that the residual concentrates on the Value-consumption
side. This clue was subsequently quantified at P2 (Section~\ref{sec:p2}), where
K-only $0.582$, V-only $1.626$, Joint $1.777$ on 40 pairs confirmed it, while
K-only being nonzero rules out Value consumption as the sole cause.

\subsection{P1 benchmark, main pair (non-equal depth, 4B$\to$1.7B)}
\label{sec:main}

\begin{table}[t]
\centering
\caption{P1 main pair (non-equal depth), three seeds, validation split (60
pairs). \textbf{Provenance-limited record}: no \texttt{458de18} rerun of this
pair exists; these values come from the superseded runs
\texttt{p1\_gpu\_main\_seed\{0,1,2\}} and are shown only to document the
broken-transfer failure boundary, not as authoritative numeric evidence. Model
pair corrected to \texttt{qwen3-4b-to-1.7b} per
\texttt{reports/runs\_index\_errata.json}. Joint $E_{\mathrm{CF}}$ [95\% CI]
reported jointly with Self and Joint dual correctness, flip consistency,
coverage, and normalized $E_{\mathrm{CF}}$.}
\label{tab:main}
\begin{tabular}{ccccccc}
\toprule
seed & Joint $E_{\mathrm{CF}}$ [95\% CI] & Self dual & Joint dual & flip & coverage & norm.\ $E_{\mathrm{CF}}$ \\
\midrule
0 & $6.044\;[4.741,7.365]$ & $0.733$ & $0.300$ & $0.490$ & $0.817$ & $0.812$ \\
1 & $5.365\;[4.156,6.663]$ & $0.650$ & $0.283$ & $0.562$ & $0.800$ & $0.827$ \\
2 & $6.080\;[4.847,7.395]$ & $0.783$ & $0.283$ & $0.542$ & $0.800$ & $0.995$ \\
\bottomrule
\end{tabular}
\end{table}

The preregistered phase-1 decision rule returns \texttt{distortion\_exists}
for all three seeds (a large gap is present and the data-quality gate passes).
However, this large gap is \emph{not} a clean utility-preserving distortion: it
coincides with a collapse of transferred dual correctness and weak direction
consistency (Section~\ref{sec:adverse}). Information attributions are mostly
negative or near zero (seed 0: mismatch $-0.488$, random $+0.229$, zero
$-0.447$; seed 1: $+1.414$, $+0.441$, $-0.040$; seed 2: $-0.248$, $-0.262$,
$-0.617$): unlike the MVP pair, the correct cache is not clearly separated
from mismatched/zero controls on this pair. No \texttt{kd\_only} baseline was
run for this pair, so we cannot distinguish ``strong baseline still distorted''
from ``transfer failed outright'' on the basis of a behavior-matched control;
the broken-transfer annotation is the basis for the failure-boundary
classification.

\subsection{P2 mechanism evidence (MVP equal-depth pair, three seeds)}
\label{sec:p2}

The mechanism stage was executed on the MVP equal-depth pair (1.7B$\to$0.6B)
with three seeds: runs \texttt{p2\_gpu\_mvp\_seed0\_delta\_rand\_audit} (the F1
pairwise-recomputation audit of seed 0), \texttt{p2\_gpu\_mvp\_seed1}, and
\texttt{p2\_gpu\_mvp\_seed2}, all at commit \texttt{6784780}, with
\texttt{zero\_re\_prefill=true} and forward counts
$\{\mathrm{capture}=0,\mathrm{oracle}=0\}$. All three runs met the E2 criteria
and yielded verdict \texttt{mechanism\_predictive}; the head-line
repair-vs-random CI was independently recomputed pairwise from the frozen
artifacts (F1 closed), and the three-seed requirement (F2) is satisfied.

\paragraph{Four arms (40 pairs per arm).} Self $0$, K-only $0.582$, V-only
$1.626$, Joint $1.777$. Only Joint has \texttt{deployable=true}; K-only and
V-only mix one mapped and one native cache half and are offline diagnostics.
V-only $\approx$ Joint quantitatively confirms the P1 direction; K-only being
nonzero means Value consumption is not the sole cause.

\paragraph{Routing/consumption decomposition.} Over 17,920 units the mean
routing norm is $24.846$ and the mean consumption norm $21.839$, with maximum
residual $2.96\times10^{-13}$ ($\approx0$ at tolerance $10^{-3}$). $d_A$ and
$d_B$ were diagnosed separately before taking differences. This is a local,
fixed-query algebraic identity, \emph{not} a global causal decomposition of
task loss.

\paragraph{Oracle held-out prediction.} For each seed the candidate unit set
$U$ was chosen on 8 training pairs (5,376 units scored) and tested on 40
held-out validation pairs with zero pair overlap, train/val prefix isolation,
and no use of the sealed test. Repair hit rates for $U$ versus random units
were $0.725/0.450$, $0.725/0.450$, and $0.750/0.400$, with paired differences
$0.275\;[0.150,0.425]$, $0.275\;[0.125,0.425]$, and $0.350\;[0.200,0.500]$
(all lower bounds $>0$) and direction consistency $0.725$/$0.725$/$0.750$.
Selected units are mostly V or both channels (no pure-K unit), concentrated at
layer 12/head 5 and layer 25/head 4. Four core units
$(12,5,\mathrm{V/both})$ and $(25,4,\mathrm{V/both})$ recur in all three seeds;
the fifth drifts marginally (seed 0: $(14,1,\mathrm{V})$, seed 1:
$(24,2,\mathrm{both})$, seed 2: $(14,1,\mathrm{both})$). We therefore describe
Value consumption as the reparable \emph{primary} locus with Key routing as an
\emph{accompanying} attention-redistribution effect --- not an exclusive or
causal claim, and not a claim of a perfectly stable unit set.

\paragraph{Interpretation, scope, and audit status.} The three E2 criteria are
\emph{not} three independent lines of evidence: repair-vs-random and direction
consistency share the unit-hit component (two statistical signals plus a
majority-threshold restatement), with the near-zero residual as a numerical
sanity check. The head-line repair-vs-random CI has been recomputed pairwise
from the frozen artifacts by the audit run
\texttt{p2\_gpu\_mvp\_seed0\_delta\_rand\_audit} (F1 closed), and the evidence
now spans three seeds (F2 closed). Hit rates ($0.725$--$0.750$) and effect
points ($0.275$--$0.350$) vary modestly across seeds, and the fifth unit drifts
marginally, so we claim core-unit stability but not per-unit invariance. All
mechanism claims are confined to the MVP equal-depth pair: the decomposition
is a local, fixed-query identity, Oracle states are offline diagnostics only,
and we make no causal or cross-model-pair assertion. Source:
\texttt{reports/runs/p2\_gpu\_mvp\_seed\{0,1,2\}/mechanism.json},
\texttt{reports/p2\_seed\_robustness.json}, and reviewer report
\texttt{reports/review/p2\_f1f2\_closure\_review.md}.

\subsection{Adverse results and boundaries}
\label{sec:adverse}

We report the following unfavorable observations explicitly, as required by the
protocol:
\begin{enumerate}
\item \textbf{Quality degradation on the main pair.} Joint dual correctness
      falls from the native $0.65$--$0.78$ to $0.28$--$0.30$ in all three seeds,
      and flip consistency is only $0.49$--$0.56$. Joint $E_{\mathrm{CF}}$
      (approx.\ $5.4$--$6.1$) is close to or above the mismatched-document and
      zero-cache controls. This is a real GPU result; we neither hide it nor
      tune around it. It means the non-equal-depth affine transfer largely
      \emph{fails}, and a high $E_{\mathrm{CF}}$ here reflects broken transfer
      rather than a well-behaved counterfactual distortion. Following the PI
      ruling (P1 Go/No-Go, \S4), this pair is recorded as a
      \textbf{transfer-failure boundary} and is \emph{not} counted as positive
      evidence for H1 (``counterfactual distortion exists''); the only current
      distortion candidate is the MVP pair, and even that candidate rests on
      three \texttt{kd\_only} seeds with an unequal training budget
      (Section~\ref{sec:pending}); it is a persistent gap across seeds, not a
      method-level improvement.
\item \textbf{Meaning of \texttt{quality\_pass} vs.\ the machine-readable
      verdict.} In this pipeline \texttt{quality\_pass=true} denotes only the
      \emph{data-quality gate} (split isolation, empty train/eval template
      overlap, annotation consistency); it contains no criterion of model
      transfer utility. Because the frozen
      \texttt{phase1\_benchmark\_decision} rule consumes
      \{distortion\_present, quality\_pass\} as inputs, the machine-readable
      verdict for the main pair is \texttt{distortion\_exists} even though the
      narrative above classifies it as broken transfer: the two differ because
      the flag and the verdict answer different questions. We therefore do not
      read the verdict alone as evidence of a clean, utility-preserving
      distortion; the metadata erratum
      (\texttt{reports/runs\_index\_errata.json}) adds an orthogonal
      \texttt{transfer\_broken} annotation without altering the frozen verdict.
\item \textbf{V-only $\approx$ Joint on MVP.} The V-only arm nearly matches
      Joint while K-only is small. E2 has since been executed (Section~\ref{sec:p2}),
      confirmed at P2 across three seeds; K-only is nonzero, so Value
      consumption is not the sole cause. The core-unit localization is stable
      across seeds, with one marginal unit drifting.
\item \textbf{Dual correctness below $0.7$ on MVP too ($0.62$--$0.67$).} The
      native Self anchor itself is imperfect on these synthetic templates; this
      is genuine model behavior not masked by any mock bias, and it would
      trigger a \texttt{tradeoff} review if carried into a method comparison.
\item \textbf{P3/P4 not run; P2 closed across three seeds.} P2 produced
      \texttt{mechanism\_predictive} in all three seeds with the pairwise CI
      recomputation audit closed (Section~\ref{sec:p2}). No \texttt{CF-RCT} training
      (though $\delta=0.5$ is locked on validation data), no sealed test, no
      generalization or cost measurement; these stages support no further
      conclusion.
\item \textbf{Huber $\delta$ locked, but no positive method direction.} The
      validation grid is complete and the PI locked $\delta=0.5$ under the
      hard-guardrail reading of ``not worse than \texttt{kd\_only}'':
      $\delta=0.5$ has dual 0.700 and flip 1.000 with ordinary QA unchanged;
      $\delta=1.0$ attains the smallest val $E_{\mathrm{CF}}$ (1.5968) but
      violates the guardrail (dual 0.650 < 0.700); $\delta=2.0$ violates it
      twice (dual 0.533, flip 0.958). All three candidate $E_{\mathrm{CF}}$
      values (2.7496 / 1.5968 / 2.7625) are above \texttt{kd\_only} (0.735):
      no $\delta$ in this grid makes CF-RCT-lite beat \texttt{kd\_only}, a real
      adverse signal for P3. The grid used 60 pairs / 3 epochs versus 120 pairs
      / 8 epochs for the \texttt{kd\_only} reference (unequal budget), and
      $\delta=0.5$ meets the guardrail exactly (``on the line''); these caveats
      are disclosed and do not justify relaxing the guardrail. P3/P5 remain
      unrun, so no method-stage conclusion is drawn.
\item \textbf{Confound caveats.} The E0 effect estimate uses multi-token
      sequence log-probabilities (scale $\sim20$) while the main endpoint uses
      2-dimensional candidate logits; because the measurement-error ceiling is
      $2.2\times10^{-16}$ this does not affect the E0 decision. The one-shot
      full-vocabulary kernel-path residual is disclosed above. The
      \texttt{model\_pair} string of the main runs is mislabeled (footnote in
      Section~\ref{sec:models}).
\end{enumerate}
Because the primary endpoint (paired $E_{\mathrm{CF}}$ vs.\ \texttt{kd\_only},
CI upper $<0$) has not been evaluated, we make \emph{no} claim of an endpoint
gain anywhere in this paper, and $P>0.05$ is nowhere interpreted as
equivalence.

\subsection{Strong baseline, $\delta$ lock, and endpoint status}
\label{sec:pending}

We summarize the available strong-baseline evidence and state explicitly which
quantities remain unevaluated.

\begin{itemize}
\item \textbf{Strong behavioral baseline (\texttt{kd\_only}): three seeds.}
      The P1 mapping arm is an affine mapper with zero trainable parameters, so
      by itself it cannot test whether a strong behavioral mapper retains the
      gap. On the MVP pair, \texttt{kd\_only} gives joint
      $E_{\mathrm{CF}}=0.735\;[0.552,0.934]$ (dual 0.700, flip 1.000, coverage
      0.800) for seed 0, $1.034\;[0.766,1.328]$ (dual 0.633, flip 1.000,
      coverage 0.817) for seed 1, and $1.539\;[1.148,1.959]$ (dual 0.683, flip
      0.979, coverage 0.800) for seed 2. All three have
      \texttt{distortion\_present=true} and \texttt{transfer\_broken=false}.
      Across the three seeds \texttt{kd\_only} narrows but does \emph{not}
      eliminate the affine gap: $E_{\mathrm{CF}}$ ranges $0.735$--$1.539$ and
      dual $0.633$--$0.700$, with every CI lower bound ($0.552$, $0.766$,
      $1.148$) above zero. Seed 0 used 120 pairs/8 epochs at commit
      \texttt{458de18}; seeds 1 and 2 used 120 pairs/8 epochs at commit
      \texttt{a04990d}, an ancestry descendant of \texttt{458de18} with
      identical run code under the round-2 audit. The unequal budget relative
      to the $\delta$ grid (60 pairs/3 epochs) remains and is disclosed, but
      the multi-seed evidence is complete. The gap being nonzero in all three
      seeds is not a finding that the strong baseline is ``close to Self'';
      equally, it does not establish a gain for CF-RCT, whose primary endpoint
      is unevaluated. No \texttt{kd\_only} run exists for the non-equal-depth
      pair.
\item \textbf{Provenance-remediation reruns (complete).} The E0 gate and the
      P1 MVP joint arm are now taken from \texttt{e0\_gpu\_mvp\_repro\_seed0}
      and \texttt{p1\_gpu\_mvp\_repro\_seed\{0,1,2\}} (all
      \texttt{git\_hash=458de18}); their values match the earlier
      \texttt{53e9bc5} runs, which are now superseded and used only as a
      cross-check. The non-equal-depth pair has no \texttt{458de18} rerun, so
      it remains a provenance-limited record.
\item \textbf{Primary endpoint: not evaluated.} The primary endpoint is the
      paired $E_{\mathrm{CF}}$ difference between a new method and
      \texttt{kd\_only} on the sealed test, with a cluster-paired 95\% CI; a
      gain may be declared only if the CI upper bound is $<0$ and the jointly
      reported metrics do not degrade. No CF-RCT model was trained, the sealed
      test was never run, and no such paired CI exists. We therefore draw
      \emph{no} conclusion about the primary endpoint, and $P>0.05$ is nowhere
      read as equivalence. The three-seed \texttt{kd\_only} evidence shows a
      persistent gap (all CI lower bounds $>0$, range $0.735$--$1.539$), so
      ``strong baseline near Self'' is not supported; however, its unequal
      budget and the absence of any CF-RCT endpoint mean no decision-tree
      branch is confirmed by this evidence alone; P3/P4 remain unrun, while P2
      is fully closed across three seeds.
\end{itemize}

% ===========================================================================
\section{Conclusion and Limitations}
\label{sec:conclusion}

We introduced an evaluation layer for counterfactual functional fidelity in
cross-model KV handoff and reported a preregistered, real-GPU benchmark.
A machine-precision measurement gate establishes that the observed effects are
not implementation artifacts. On an equal-depth pair, transfer exhibits a
counterfactual gap whose CI excludes zero while ordinary direction
consistency is largely preserved. On a non-equal-depth pair, the same gap is
large but accompanied by a severe loss of answer correctness, delimiting a
failure boundary for the current affine transfer.

The evidence supports contribution (1) and contribution (2), but not
contribution (3). Contribution (1) is limited to the affine/ridge
mapper tested here; the three-seed \texttt{kd\_only} evidence narrows that gap
without eliminating it, and the other strong baselines
(\texttt{consumer}, \texttt{score}, ridge) are not evaluated. Contribution (2)
rests on three-seed, MVP-pair held-out predictive evidence localizing the gap
to receiver-side Value-consumption pathways with accompanying Key routing,
reported without causal or cross-model-pair claims (Section~\ref{sec:p2}).
There is \emph{no} trained
method evaluated on the paired endpoint; the CF-RCT design is presented
strictly as pending validation. The non-equal-depth results are a transfer-failure
boundary and are not counted as positive evidence for H1. All experiments use
synthetic templates, so the applicability of these observations to other
domains or model families is not established. The main numeric tables come
from the reproducible \texttt{458de18} runs described in
Section~\ref{sec:models}; the non-equal-depth pair has no such rerun and is
retained only as a provenance-limited failure-boundary record. Across three
seeds \texttt{kd\_only} narrows the affine gap without eliminating it (all CI
lower bounds above zero); the unequal budget remains a caveat and the locked
$\delta=0.5$ grid gives no positive method direction. Following the
decision tree
(protocol~\S8), this work should
be written and positioned as a \emph{boundary / empirical study}, not as a
high-level method paper. Future work would need to (i)~test the mechanism
localization on non-equal-depth pairs (E5/P4) before any cross-model-pair claim;
(ii)~train CF-RCT-lite/mech, lock $\delta$ on validation, and evaluate the
paired endpoint against \texttt{kd\_only} with $\ge3$ seeds and Holm; and
(iii)~run generalization and cost stages before any claim of deployability or
speedup. Until then, contribution (3) is withheld.

% ===========================================================================
\section*{Appendix A: Figures, tables, and submission materials}

Figure~\ref{fig:f1}--\ref{fig:f3} are textual placeholders; the complete figure
plan (Figures~1--7) and Tables~1--3, a cover-letter outline, and reviewer
response templates are maintained in \texttt{papers/readkv\_v2\_draft/notes.md}.
All figures and tables must be generated directly from the frozen result files
without recomputing statistics. Figures~4--7 (native vs.\ transferred response
scatter, layer/head error heatmap with Oracle repair, method ablation,
generalization and cost--Pareto) are deferred until their underlying stages are
actually run; they must not be fabricated. Figure~5 material (layer/head error
heatmap with Oracle repair) is available from the P2 run
the P2 runs (\texttt{p2\_gpu\_mvp\_seed\{0,1,2\}} and the audit run
\texttt{p2\_gpu\_mvp\_seed0\_delta\_rand\_audit}; \texttt{mechanism.json},
\texttt{mechanism\_units.jsonl}, \texttt{decomposition.jsonl}) and may be drawn
from those frozen artifacts; Figures~6--7 await P3/P4.

\bibliographystyle{plain}
\begin{thebibliography}{9}
\bibitem{kvlingo} KV-Lingo. \emph{Citation placeholder---verify formal version,
date, and experimental protocol before submission (protocol \S9).}
\bibitem{cachelink} CacheLink. \emph{Citation placeholder---verify formal
version, date, and experimental protocol; report deviations if no runnable
implementation exists.}
\end{thebibliography}

\end{document}
